Les assistants de code changent en profondeur la façon dont on écrit un texte
mathématique. Il y a encore peu, un modèle de langage produisait au mieux un brouillon à
reprendre ligne à ligne. Aujourd'hui, un assistant qui travaille dans un dépôt lit des
articles, transcrit des notes manuscrites, écrit du \LaTeX{} qui compile, énonce des
théorèmes dans un assistant de preuve et les démontre, cherche dans la littérature et
rédige une première version. Donnez-lui un rapport de stage et quelques articles publics
sur le sujet, avec un dépôt qui cadre bien ce qu'on attend : en une dizaine de minutes, il
en tire une leçon exploitable. De nouvelles techniques d'écriture automatique apparaissent,
et elles demandent une méthode.

Cette méthode, je l'ai cherchée en mathématicien qui a quitté la recherche après son master
pour faire carrière dans l'ingénierie des données. J'en ai gardé deux choses : un goût pour
les mathématiques, et l'habitude des outils qui rendent un programme fiable. La
programmation fonctionnelle m'avait appris qu'on peut écrire du code « comme en maths », de
façon rigoureuse ; les assistants de preuve poussent l'idée jusqu'au bout. Pendant quelques
mois, et surtout pendant des week-ends, j'ai mené avec un assistant, et avec de nombreuses
sessions lancées en parallèle dans le cloud, un projet qui m'aurait autrefois demandé une
thèse : transcrire plusieurs milliers de pages de notes mathématiques manuscrites, les
lire, chercher dans la littérature ce qu'elles contenaient de connu, et faire démontrer
par la machine les énoncés qu'on ne trouvait nulle part. Une bonne partie portait sur un
domaine que je n'ai jamais étudié.

C'est ce dernier point qui m'a convaincu d'écrire ce livre. Si un non-spécialiste peut
compiler de tels recueils de résultats, la question n'est plus seulement mathématique ;
elle est de méthode. Un assistant écrit un texte très convaincant, et vrai par construction
la plupart du temps. C'est déconcertant, et cela demande de la vigilance : une erreur s'y
glisse avec la même assurance que le reste. Il faut donc savoir à chaque instant ce qui a été vérifié,
par qui, et comment. Les revues commencent à fixer des règles (certaines refusent les
textes écrits avec l'intelligence artificielle), les sites d'entraide aussi, et les
enjeux pour la formation et l'emploi des chercheurs sont réels. Il manque une sorte de code
de la route pour les chercheurs qui utilisent ces outils et veulent partager l'espace public
du débat scientifique sans le dégrader. Ce livre n'a pas cette prétention, mais il en
propose quelques articles.

Il s'adresse à une mathématicienne ou à un mathématicien qui mène un projet trop gros pour
tenir dans sa tête : l'édition d'un corpus de notes, la formalisation d'une théorie dans un
assistant de preuve, un livre écrit à plusieurs, un séminaire dont on veut garder les
traces. Il décrit les outils et les pratiques du \emph{software craftsmanship} qui rendent
un tel projet tenable, et la façon dont on s'en sert quand une partie du travail est
confiée à un assistant. Le projet qui m'a servi n'en est pas le sujet : il fournit des
exemples, réduits à la leçon qu'ils portent, et quelques chiffres pour donner l'échelle.
J'ai appris la plupart de ces règles en commettant l'erreur qu'elles évitent.

On ne suppose aucune formation en informatique. Il faut savoir ouvrir un terminal et taper
une commande, et accepter de le faire souvent. Chaque chapitre part d'un problème que pose
le travail mathématique (retrouver pourquoi un énoncé a changé, savoir ce qui a été vérifié
et par qui, ne pas publier ce qu'on n'a pas le droit de publier) et montre l'outil qui le
règle, puis la discipline sans laquelle l'outil ne sert à rien.

Car l'outil, sans discipline, ne sert pas à grand-chose. Un dépôt \texttt{git} où chaque
commit s'intitule « modifications » ne garde rien de ce qu'on voudra savoir dans trois
mois. Une liste d'issues que personne ne ferme tourne au grenier. Et un assistant sans
règles écrites produit beaucoup, très vite, sans qu'on sache ensuite ce qui a été relu. La
plupart des pages qui suivent décrivent donc des règles, chacune avec sa raison.

Ces pratiques seront sans doute courantes d'ici un an ou deux, et d'autres les auront
décrites mieux que moi. Je les écris maintenant parce qu'elles servent maintenant.

Le livre se lit dans l'ordre, mais chaque chapitre tient seul. L'annexe rassemble les
commandes citées.

\begin{noteblock}
Les exemples de commandes sont donnés pour un terminal Unix (macOS ou Linux). Les noms
d'outils commerciaux sont ceux qu'on a utilisés ; d'autres font le même travail, et les
règles valent pour eux aussi.
\end{noteblock}
